% =====================================================================
% Keccak-f[1600] in Microcode — restructured into three sections:
%   1. Designing the Microcode Kernel   (\label{sec:keccak-design})
%   2. Implementation                   (\label{sec:keccak-impl})
%   3. Evaluation                       (\label{sec:keccak-eval})
%
% Worked example instantiating the methodology of Section 3
% (\Cref{sec:ucode-layer}) on a control-flow-heavy primitive. Numbers and
% design rationale are taken from docs/keccak-microcode-rundown.html and the
% generated patch (keccak_perm_body.h / keccak_gen.py).
%
% Assumes the Background has recalled the Keccak-f[1600] round (theta, rho,
% pi, chi, iota); we restate only what the mapping needs.
% =====================================================================

% =====================================================================
\section{Designing the Microcode Kernel}
\label{sec:keccak-design}
% =====================================================================

The first worked example instantiates the methodology of
\Cref{sec:ucode-layer} on the Keccak-$f[1600]$ permutation. Keccak is the most
demanding test of the layer in this work: in contrast to the straight-line
field-arithmetic primitives it is a 24-round loop with a per-round constant
table, and it therefore exercises in-place looping, indexed addressing, and
conditional control flow in addition to dense arithmetic packing. It is also
the example in which the structural advantage of microcode is most readily
stated, and, as we shall see, the example in which the limits of that advantage
are most clearly exposed. This section sets out the design decisions; the
concrete mechanics that realise them are deferred to \Cref{sec:keccak-impl}.

\paragraph{The structural opportunity: residency across all rounds.}
A Keccak state consists of 25 lanes of 64 bits. Native x86-64 exposes
approximately 14 usable general-purpose registers, with the consequence that
hand-written assembly cannot hold the entire state in registers; it spills to
the first-level cache continuously, and the fastest scalar implementations
recover residency only by fully unrolling all 24 rounds. The
micro-architectural register file is 32 entries wide, comprising 16
general-purpose and 16 temporary registers, which is sufficient to hold the
complete 25-lane state resident not merely within a single round but across all
24. Our first and governing design decision is to rest the entire kernel on
this property: the 25 lanes are loaded once, the 24 rounds are evaluated with
the whole state resident in registers, and the lanes are stored once at the
end. Residency is obtained without unrolling, which the 128-triad capacity of
\Cref{sec:budget} would not in any case permit, and the resulting register
headroom is then spent on per-round transformations that register-starved
native code is unable to perform.

\paragraph{One firing for the whole permutation.}
The second decision concerns where the round boundary lives. Rather than fire
the patch once per round, which would force the state out to memory and back at
every boundary and pay the hook cost 24 times, we evaluate the entire
permutation within a \emph{single} hooked \texttt{vmwrite} firing and loop the
round body internally by means of a backward branch. Because control never
leaves the patch, the temporary registers persist for the lifetime of the
permutation, and the state therefore stays resident across all 24 rounds with
no boundary traffic. Only a prologue and an epilogue, executed once each, move
the state between registers and memory. This choice is what makes the residency
of the previous paragraph realisable in practice.

\paragraph{Why a single looped body is correct: the structure of $\pi$.}
The natural objection to looping a Keccak round is that $\pi$ relabels lane
positions, so that the lane held in a given register would appear to drift from
one round to the next, which would seem to demand either per-round register
renaming or two alternating body copies. The objection is unfounded, and the
reason is structural. The permutation $\pi$ on the 25 lanes consists of one
fixed point, namely lane~0, together with a single 24-cycle. If $\pi$ is
applied \emph{in place} along that cycle --- that is, by walking the cycle and
moving each element into its destination register while folding in the
exclusive-or of $\theta$ and the rotation of $\rho$ --- then every lane is
returned to its canonical register by the end of the round. Round $r{+}1$
consequently observes exactly the register layout that round $r$ observed, so
that a single identical body may be looped 24 times with no renaming and no
two-body alternation scheme. This observation is the linchpin of the design: it
is what allows the static footprint to be one round body rather than 24, and
hence what allows residency and the 128-triad budget to coexist.

\paragraph{Borrowing a 32nd register.}
Holding the 25 state lanes, the five mixing lanes $D$ of $\theta$, the buffer
base, and a working temporary simultaneously requires one register more than
the architecturally nameable file affords once the base pointer is reserved.
Our third decision is therefore to borrow \texttt{RSP} as a 32nd data register,
saving it to the buffer in the prologue and restoring it in the epilogue. This
is safe because, as established in \Cref{sec:findings}, the patch never touches
the stack. The substitution is not cosmetic: it is precisely what permits the
five $D$ values to remain resident in registers throughout a round instead of
being re-read from memory once per lane, which would otherwise cost 25
additional memory triads per round.

\paragraph{Spending the headroom.}
The final design decisions concern how the freed registers are used, each being
a transformation that native code cannot make because it lacks the spare
registers to do so. The state is kept fully resident, so that input and output
are paid once per permutation rather than once per round. The five $D$ values
are kept resident, as just described. And $\chi$, which would otherwise force
two auxiliary moves per row to preserve values it is about to overwrite, is
restructured to use the now-dead $D$ registers as scratch and so requires no
auxiliary moves at all. To these we add two latency-shortening choices on the
round's critical path: a balanced parity tree for $\theta$, and a loop counter
that stores the round constant's byte offset directly so that the constant
fetch needs no scaling arithmetic. The mechanics of all five appear in
\Cref{sec:keccak-impl}, and their measured contributions in
\Cref{sec:keccak-eval}. One further decision is forced rather than chosen: the
loop counter is advanced \emph{upward}, because the decrement-style operations
that a count-down would require do not behave correctly on this part, as
recorded in \Cref{sec:findings}.

% =====================================================================
\section{Implementation}
\label{sec:keccak-impl}
% =====================================================================

We now give the concrete realisation of the design of \Cref{sec:keccak-design}:
the register and buffer layout, the triad-level form of each round stage, the
generator and packer that produce the patch, and the verification that
accompanies it.

% ---------------------------------------------------------------------
\subsection{State Layout and Residency}
\label{sec:keccak-layout}

\paragraph{Register allocation.}
The 25 lanes, indexed by $i = 5y + x$, are assigned to a fixed set of 25
register names, namely the 13 general-purpose registers
$\texttt{RDI},\texttt{RSI},\texttt{RBX},\texttt{RDX},\texttt{RBP},
\texttt{R8}\ldots\texttt{R15}$ followed by the 12 temporaries
$\texttt{TMP0}\ldots\texttt{TMP11}$. The register \texttt{RCX} holds the buffer
base and is never repurposed, while the remaining five registers, namely
$\texttt{RAX}$ and $\texttt{TMP12}\ldots\texttt{TMP15}$, serve as scratch: they
hold the five resident $D$ values of $\theta$ during a round and the temporaries
of $\chi$ once those values are no longer needed. With \texttt{RSP} borrowed as
the 32nd data register (\Cref{sec:keccak-design}) the file is exactly filled,
\texttt{RSP} serving as the working register for the parity tree, the $\rho$
rotation, and the $\pi$ cycle walk.

\paragraph{Buffer organisation.}
The operation is backed by a single contiguous array of 56 lanes, with the base
pointer centered at \texttt{RCX}~$=$~\texttt{\&buf[16]} so that every
displacement fits within the 8-bit signed offset field discussed in
\Cref{sec:findings}, as summarised in \Cref{tab:keccak-buffer}. Only the 25
state lanes are accessed by the prologue and epilogue. The column-parity
scratch is spilled and re-read within a round, whereas the round-constant table
lies beyond the reach of the immediate offset and is therefore addressed
through the register-indexed load form. The loop counter is stored not as a
round number but as the byte offset of the current constant, so that the
constant fetch requires no scaling arithmetic on its critical path.

\begin{table}[t]
\centering
\caption{Operand buffer of 56 lanes, with the base register centered at
\texttt{RCX}$=$\texttt{\&buf[16]}.}
\label{tab:keccak-buffer}
\small
\begin{tabular}{@{}llll@{}}
\toprule
lane & offset(\texttt{RCX}) & name & contents \\
\midrule
0--24  & $-128\ldots{+}64$ & state    & resident in 13 GPR $+$ 12 TMP during a round \\
25--29 & $+72\ldots{+}104$ & DSCR     & $\theta$-$C$ parity scratch (spilled, re-read) \\
30     & $+112$            & COUNTER  & RC byte-index, initialised to $128$ and advanced by $8$ each round \\
31     & $+120$            & RSPSAVE  & saved \texttt{RSP} (borrowed as a data register) \\
32--55 & (indexed)         & RC table & 24 round constants, addressed only by register index \\
\bottomrule
\end{tabular}
\end{table}

% ---------------------------------------------------------------------
\subsection{The Looped Round Body}
\label{sec:keccak-round}

The permutation is the prologue, a body looped 24 times by a backward branch,
and the epilogue. The prologue loads the 25 lanes using one single-load triad
per lane, since at most one memory operation may share a triad
(\Cref{sec:keccak-packing}); the epilogue stores them symmetrically. We
describe the body in order of execution.

\paragraph{$\theta$ and the residency of $D$.}
The column parity $C_x=\bigoplus_y A[x{+}5y]$ stands at the head of every
round's dependency chain, so we shorten its latency. In place of the reference
implementation's depth-four left-associated chain we emit a balanced tree of
depth three, using \texttt{RSP} as a second accumulator so that the two halves
are computed in parallel, and the packer places the dependent exclusive-or
operations in successive slots of a single triad, as shown in
\Cref{lst:keccak-theta}. The five values $C_x$ are then spilled to the buffer,
because the registers are required immediately afterwards, and the mixing lanes
$D_x = C_{x-1}\oplus \mathrm{rol}(C_{x+1},1)$ are computed into the five scratch
registers and kept resident for the remainder of the round, eliminating the 25
per-lane reloads of $D$ that earlier versions of the kernel incurred.

\begin{lstlisting}[language=C,caption={The depth-three parity tree of $\theta$,
in which a single column's tree is packed into the slots of one triad while the
spill of the previous column is interleaved.},label={lst:keccak-theta},basicstyle=\ttfamily\footnotesize]
{ XOR(RAX,RDI,R8),   XOR(RSP,R13,TMP2),  XOR(RAX,RAX,RSP),  NOP_SEQW }
{ XOR(RAX,RAX,TMP7), STAD(RAX,RCX,72),   XOR(RAX,RSI,R9),   NOP_SEQW }
\end{lstlisting}

\paragraph{The fused $\theta$-apply, $\rho$, and $\pi$ stage.}
By the argument of \Cref{sec:keccak-design}, $\pi$ is one fixed point plus a
single 24-cycle, so applying it in place along that cycle returns every lane to
its canonical register. Concretely, we walk the cycle and move each element into
its destination register while folding in the exclusive-or of $\theta$ and the
rotation of $\rho$, reading $D$ directly from registers with no memory access.
The fusion costs one exclusive-or and one rotation per lane, and \texttt{RSP}
holds the one element displaced at the head of the cycle.

\paragraph{The $\chi$ stage without auxiliary moves.}
The step $\chi$ computes, for each row, $A_x \mathbin{\oplus}\!= (\lnot
A_{x+1})\land A_{x+2}$, which corresponds exactly to the \texttt{NOTAND}
primitive. Performed naively in place, the overwriting of $A_x$ would destroy a
value the subsequent lanes still require, which forces two auxiliary moves per
row, or ten per round, in register-starved native code. Because the five $D$
registers are now dead, we have the headroom to compute all five \texttt{NOTAND}
terms of a row first, placing them in scratch, and then to apply the five
in-place exclusive-or operations, as shown in \Cref{lst:keccak-chi}, so that no
auxiliary moves are required at all. This is the single largest measured
improvement (\Cref{sec:keccak-eval}), and it is a direct consequence of
residency.

\begin{lstlisting}[language=C,caption={The move-free $\chi$ stage, in which the
five \texttt{NOTAND} terms of a row are computed into the now-dead
$D$-registers before the five in-place exclusive-or operations are
applied.},label={lst:keccak-chi},basicstyle=\ttfamily\footnotesize]
{ NOTAND(RAX,RSI,RBX), NOTAND(TMP12,RBX,RDX), NOTAND(TMP13,RDX,RBP), NOP_SEQW }
{ NOTAND(TMP14,RBP,RDI),NOTAND(TMP15,RDI,RSI),XOR(RDI,RDI,RAX),      NOP_SEQW }
{ XOR(RSI,RSI,TMP12),   XOR(RBX,RBX,TMP13),   XOR(RDX,RDX,TMP14),    NOP_SEQW }
\end{lstlisting}

\paragraph{$\iota$ and loop control.}
Because the counter lane stores the byte offset of the current round constant
directly, $\iota$ reduces to a load of the index, an indexed load of the
constant, and a single exclusive-or into lane~0, with no scaling arithmetic on
the path from $\iota$ to lane~0. The counter is then advanced by 8 and persisted
in the spare slots of the same triad, after which a single forced triad performs
the loop test and the backward branch. In that triad an exclusive-or of the
index against the end marker sets the architectural zero flag, and a conditional
branch within the same triad returns control to the top of the body for as long
as the index differs from $320$, which lies just past the last constant, thereby
producing exactly 24 iterations, as shown in \Cref{lst:keccak-loop}.

\begin{lstlisting}[language=C,caption={Loop control, in which the flag-setting
exclusive-or and the backward conditional branch must occupy the same
triad.},label={lst:keccak-loop},basicstyle=\ttfamily\footnotesize]
{ XOR(TMP12,TMP14,320), NOP, UJMPCC_CONDNZ(TMP12, 0x7c68 /*loop_top*/), NOP_SEQW }
\end{lstlisting}

% ---------------------------------------------------------------------
\subsection{Generation, Packing, and the Triad Budget}
\label{sec:keccak-packing}

\paragraph{No hand-written triads.}
In accordance with \Cref{sec:triads-packing}, no triad is written by hand. The
generator emits the round as an ordered op-list over real register names, and a
single greedy left-to-right packer lowers that list to triads, closing the
current triad whenever it would exceed three slots, whenever a second memory
operation would enter it, whenever an operation would consume a register loaded
earlier in the same triad, and around the forced loop-test triad. The packer
therefore enforces the hardware scheduling rules mechanically rather than
leaving them to the programmer.

\paragraph{The one-memory-operation rule.}
The packer admits at most one memory operation per triad, which is why the
prologue and epilogue are 25 single-lane triads each. A targeted probe refined
the cause of this rule. A single triad bearing two memory operations executes
correctly in every combination of loads and stores we tested in isolation;
however, a kernel in which many such triads are packed consecutively --- as
arises naturally when the 25-lane prologue and epilogue are compressed two lanes
to a triad --- faults the machine. The binding constraint is thus the sustained
throughput of the memory cluster rather than the legality of an individual
triad, and we retain the one-operation-per-triad rule in production.

\paragraph{Budget.}
The complete permutation compiles to 108 triads, which lies comfortably within
the 128-triad patch RAM. These comprise a 26-triad prologue, consisting of one
\texttt{RSP} save and 25 single-lane loads, a 56-triad body that includes the
constant lookup and the loop test, and a 26-triad epilogue, consisting of 25
stores and the \texttt{RSP} restore. The body begins at \texttt{U7c68}, which is
the branch target. Because the body is looped rather than unrolled, the static
footprint is independent of the number of rounds, whereas the dynamic count ---
the number of triads actually executed per permutation --- is $26 +
24\times56 + 26 = 1396$. The 20 triads of headroom beneath the cap are genuine
slack, but, as \Cref{sec:keccak-eval} explains, they cannot be spent on
unrolling, because even a single additional round body would not fit alongside
the existing one.

% ---------------------------------------------------------------------
\subsection{Verification}
\label{sec:keccak-verify}

Before any triads are emitted, the generator executes the op-list on a
register-and-memory simulator that models the 8-bit signed-offset limit, so that
an out-of-range displacement manifests as an assertion failure in software
rather than as a silent corruption of state on hardware, and asserts equality
with a reference Keccak over 64 random states across all 24 rounds. On hardware
the installed patch is then checked in two ways: first lane by lane against the
same reference, and second, more stringently, against the published $f(0)$ and
$f(f(0))$ known-answer vectors of the Keccak team, so that a correlated error
present in both the reference and the patch cannot escape detection. The
all-zero anchor reproduces the canonical value
$\texttt{lane0}=\texttt{0xF1258F7940E1DDE7}$.

% =====================================================================
\section{Evaluation}
\label{sec:keccak-eval}
% =====================================================================

\paragraph{Measurement methodology.}
Timing on this part is subject to a pitfall that invalidates naive comparisons.
The \texttt{rdtsc} instruction counts at a fixed reference rate of approximately
$1.1$\,GHz, whereas the core bursts to approximately $2.4$\,GHz under load, with
the consequence that the same work is reported as roughly $2.2$ times more or
fewer ticks depending on the power state during measurement. Comparing a fresh
measurement against a baseline captured in a separate run is therefore
meaningless. We instead time the microcode and all of the scalar contenders back
to back within a single process, after a warm-up phase has driven the core to a
steady frequency, so that both observe the same clock and the ratio between them
is frequency-invariant; for absolute cycle counts we additionally pin the core
to its base frequency so that the reported ticks approximate true cycles. We
compare against the genuinely fastest scalar baselines from SUPERCOP, swept over
24 compiler-and-optimisation configurations with the best result selected for
each contender, in the same manner as the SUPERCOP autotuner.

\paragraph{Result.}
Under base-frequency, same-process measurement the looped permutation completes
in 1910 cycles per permutation, as against 2045 cycles for the fastest scalar
contender, \texttt{opt64lcu24}, and 2061 cycles for the hand-tuned
\texttt{x86\_64\_asm}, as reported in \Cref{tab:keccak-results}. This
corresponds to a ratio of $0.934$, an improvement of approximately $7\%$, and
the microcode is the fastest of all contenders. The ratio is the robust
quantity, since the same harness measured at burst frequency reports
approximately 874 cycles against 942, the identical ratio of $0.93$ scaled by
the $2.4/1.1$ power-state factor.

\begin{table}[t]
\centering
\caption{Best result for each contender over 24
compiler-and-optimisation configurations, measured base-pinned and
same-process, where a smaller value is better.}
\label{tab:keccak-results}
\small
\begin{tabular}{@{}lrrl@{}}
\toprule
contender & min (cyc/perm) & median & best config \\
\midrule
\textbf{microcode (looped)}      & \textbf{1910} & 1916  & gcc-13 -Os \\
\texttt{opt64lcu24}              & 2045          & 2054  & gcc-11 -Os \\
\texttt{x86\_64\_asm}            & 2061          & 2069  & clang-14 -O2 \\
\texttt{opt64lcu24shld}          & 10053         & 10064 & clang-17 -Os \\
\texttt{x86\_64\_shld}           & 10086         & 10097 & gcc-13 -O2 \\
\bottomrule
\end{tabular}
\end{table}

\paragraph{The source and the bound of the improvement.}
The margin is modest by design, and the latency model of \Cref{sec:findings}
explains why. A single isolated firing sustains $0.49$ cycles per triad, but
this is a throughput figure obtained by overlapping independent firings, whereas
the looped permutation is a strict dependency chain, in which round $N$ requires
round $N{-}1$, with a backward branch on each round and no overlap between
firings, so that it pays the true latency of approximately $1.37$ cycles per
triad and the 1396 dynamic triads accordingly cost approximately 1910 cycles.
The microcode is therefore latency-bound at roughly one triad per cycle, whereas
the hand-tuned scalar assembly extracts approximately three operations per cycle
from the rotate and \texttt{ANDN} instruction-level parallelism of Goldmont. We
cannot close this gap with additional parallelism, because the chain
$\theta{\to}\rho{\to}\pi{\to}\chi$ is inherently sequential and the 128-triad
cap precludes the unrolling that would expose independent work. The entire
advantage is consequently the per-round work that native code cannot perform:
full register residency, with the state input and output paid once per
permutation rather than once per round; the borrowed 32nd register that keeps
$D$ resident; and the move-free $\chi$ stage. \Cref{tab:keccak-opt} shows the
improvement assembled from these structural contributions, each measured back to
back.

\begin{table}[t]
\centering
\caption{The improvement assembled from successive structural optimisations,
expressed as the ratio to the naive residency-only baseline.}
\label{tab:keccak-opt}
\small
\begin{tabular}{@{}clr@{}}
\toprule
step & change & ratio \\
\midrule
0 & per-lane $D$ in buffer, with a naive $\chi$ using auxiliary moves & $1.00\times$ \\
1 & $D$ resident in registers via the borrowed \texttt{RSP}, with no base rebuild & $0.99\times$ \\
2 & move-free $\chi$ using the dead $D$-registers as scratch & $0.95\times$ \\
3 & depth-three $\theta$ tree and byte-indexed RC, shortening the critical paths & $0.93\times$ \\
\bottomrule
\end{tabular}
\end{table}

In summary, a complete and known-answer-verified 24-round Keccak-$f[1600]$ runs
entirely in Goldmont microcode, looped within a single \texttt{vmwrite} firing,
in 108 triads, and outperforms the fastest hand-tuned scalar assembly on this
microarchitecture by approximately $7\%$. The result is a clear demonstration of
the character of the layer established in \Cref{sec:ucode-layer}: its advantage
is structural, deriving from the elimination of the spills and data movement
that the instruction set architecture cannot avoid, rather than from a higher
peak throughput, and it is greatest precisely where native code is most
constrained by the shortage of registers.
